% year: תשפ"ה
% type: בוחן
% semester: א
% moed: לדוגמה
% official_solution: yes
% source: exams/מבחנים/תשפו/midterm_example 2025.pdf
% part: ב
% number: 5
% points: 10
% categories: functions, multiple-choice

\begin{question}
תהי
\[
f(x)=\begin{cases}
3x-1, & x<0,\\
x^2, & x>0
\end{cases}
\]
אילו מהטענות הבאות נכונה?
\begin{enumerate}
  \item הפונקציה אינה מונוטונית.
  \item הפונקציה מונוטונית יורדת ממש.
  \item הפונקציה מונוטונית עולה ממש.
  \item הפונקציה מונוטונית אך לא ממש.
\end{enumerate}
\end{question}

\begin{hint}
בדקו כל חלק בנפרד, ואז השוו ערך כלשהו משמאל ל-0 לערך כלשהו מימין ל-0.
\end{hint}

\begin{solution}
\textbf{התשובה הנכונה: (ג)}.

(שימו לב: כפי שהפונקציה כתובה, היא אינה מוגדרת ב-$x=0$; בפתרון המצורף נכתב $x\le0$ בחלק הראשון. התשובה זהה בשני המקרים.)

יהיו $x_1<x_2$ בתחום ההגדרה. נבחין בשלושה מקרים:
\begin{itemize}
  \item $x_1<x_2<0$: $f(x_2)-f(x_1)=3(x_2-x_1)>0$.
  \item $0<x_1<x_2$: $f(x_2)-f(x_1)=x_2^2-x_1^2=(x_2-x_1)(x_2+x_1)>0$, כי שני הגורמים חיוביים.
  \item $x_1<0<x_2$: $f(x_1)=3x_1-1<-1<0<x_2^2=f(x_2)$.
\end{itemize}
(אם מגדירים $f(0)=3\cdot0-1=-1$, גם המקרים $x_1<0=x_2$ ו-$x_1=0<x_2$ מקיימים $f(x_1)<f(x_2)$ באותו אופן.)
בכל המקרים $f(x_1)<f(x_2)$, ולכן $f$ עולה ממש.

\textbf{למה האחרות שגויות:} (א) — הראינו ש-$f$ מונוטונית. (ב) — למשל $f(1)=1>f(-1)=-4$, ולכן אינה יורדת. (ד) — "מונוטונית אך לא ממש" אומר שקיימים $x_1<x_2$ עם $f(x_1)=f(x_2)$, וזה לא קורה כי $f$ עולה ממש.
\end{solution}
